% !TEX root = ../main.tex
\section{Lenguaje y organización de la solución}\label{sec:solucion}

\subsection{Lenguaje declarado}

El proyecto se desarrolla en \Lenguaje\ sobre la plataforma \Plataforma.

\subsection{Estructura de la solución}

La solución se divide en proyectos cuyas dependencias apuntan en una sola dirección, hacia el dominio. De este modo la lógica del juego no conoce la consola ni ningún otro medio de entrada y salida. La Tabla~\ref{tab:proyectos} describe cada proyecto.

\begin{table}[H]
\caption{Proyectos que conforman la solución}\label{tab:proyectos}
\tablacompacta
\begin{tabularx}{\textwidth}{@{}L{3.7cm}Y L{3.3cm}@{}}
  \toprule
  \cab{Proyecto} & \cab{Responsabilidad} & \cab{Depende de} \\ \midrule
  \texttt{Domino.Dominio} & Reglas y conceptos del juego: partida, mesa, fichas, fichas para robar, jugadores y turnos. & Ninguno \\ \filasep
  \texttt{Domino.Persistencia} & Guardado y carga de partidas en archivos JSON, detrás de \texttt{IRepositorioPartidas}. & \texttt{Domino.Dominio} \\ \filasep
  \texttt{Domino.Consola} & Interfaz de texto: presenta el estado del juego y traduce las acciones del jugador. & \texttt{Domino.Dominio}, \texttt{Domino.Persistencia} \\ \filasep
  \texttt{Domino.Simulacion} & Ejecución de partidas completas sin interfaz, con semilla fija, para la medición. & \texttt{Domino.Dominio}, \texttt{Domino.Persistencia} \\
  \bottomrule
\end{tabularx}
\end{table}

\begin{nota}
\textbf{Restricción de independencia del dominio.} El proyecto \texttt{Domino.Dominio} no referencia a \texttt{System.Console} ni a ningún otro proyecto de la solución. La consola puede reemplazarse por otra interfaz sin modificar el dominio, y \texttt{Domino.Simulacion} lo evidencia al ejecutar partidas completas utilizando únicamente el dominio.
\end{nota}

\subsection{Dependencias externas}

El control del tiempo por turno y la persistencia de la partida utilizan componentes de la plataforma, pero el dominio los consume a través de abstracciones propias (\texttt{IReloj} e \texttt{IRepositorioPartidas}). De este modo la simulación emplea un reloj controlado, que no espera segundos reales, y el medio de guardado puede cambiarse sin afectar las reglas. La persistencia reside en un proyecto propio porque la utilizan tanto la consola como la simulación, que mide el guardado y la carga.

\subsection{Estado del código}

El repositorio (\Repositorio), con la etiqueta \EtiquetaRepo, contiene en la carpeta \texttt{Documentación} este documento con sus diagramas y, en la carpeta \texttt{Código}, el esqueleto inicial de la solución \texttt{Domino.slnx}. El esqueleto incluye los proyectos \texttt{Domino.Dominio} y \texttt{Domino.Consola}, compila sin advertencias y se ejecuta desde la consola.

En esta entrega el dominio contiene las clases \texttt{ControladorJuego}, \texttt{Partida}, \texttt{Jugador}, \texttt{JugadorHumano}, \texttt{JugadorCpu}, \texttt{Mano}, \texttt{Ficha} y \texttt{ConfiguracionPartida}, junto con las enumeraciones \texttt{EstadoPartida}, \texttt{TipoDomino} y \texttt{NivelDificultad}, con los mismos nombres del diagrama de clases. La consola ejecuta la parte central de CU-01 Configurar partida: solicita el alias, el tipo de dominó, el número de CPU y el nivel de cada una, valida el alias y el número de CPU (RN-02) y crea la partida en estado Preparada con el jugador humano en la posición~0 y las CPU a continuación. La elección del límite de tiempo, el reparto de las fichas, el ciclo de turnos y los proyectos \texttt{Domino.Persistencia} y \texttt{Domino.Simulacion} se incorporan en las entregas siguientes, a medida que se implementan las clases restantes del diagrama.
